\begin{Exercise}[label={ex:cifrario}]
    Si scriva un programma che decifra o cifra un file di testo applicando uno scorrimento ai caratteri alfabetici.
    Il programma deve accettare i parametri \emph{in qualsiasi ordine} tramite dei flag da linea di comando:
    \begin{itemize}
        \item \texttt{-i <inputfile>} : specifica il file di input
        \item \texttt{-o <outputfile>} : specifica il file di output
        \item \texttt{-k <shift>} : specifica la chiave
        \item \texttt{-d} : decifra il file (se non presente, cifra)
    \end{itemize}

    \textbf{Esempio:}
    \begin{lstlisting}[style=console]
        # Cifra il file msg.txt con chiave 3 e salva il risultato in cif.txt
        ./programma -i msg.txt -o cif.txt -k 3
        
        # Decifra il file cif.txt con chiave 5 e salva il risultato in msg.txt
        ./programma -k 5 -d -o msg.txt -i cif.txt
    \end{lstlisting}
\end{Exercise}

\begin{Answer}[ref={ex:cifrario}]
  	\lstinputlisting[style=cstylesmall]{src/cifrario.c}
\end{Answer}